\documentclass[11pt,a4paper]{article}
\usepackage[margin=1in]{geometry}
\usepackage{amsmath,amssymb}
\usepackage{cite}

\title{Hierarchical Cortical Circuit Dynamics Underlying Working Memory Consolidation}
\author{Aria Thorne, David Sterling, and Kenji Takahashi}
\date{February 2026}

\begin{document}
\maketitle

\begin{abstract}
Working memory maintenance requires coordinated recurrent excitation within prefrontal microcircuits. Here, using two-photon calcium imaging and high-density neuropixels probes in behaving primates, we identify a distinct subpopulation of deep-layer pyramidal neurons whose persistent firing directly encodes task-relevant mnemonic representations \cite{smith2024, brown2023}. Optogenetic suppression of layer 5 projections selectively abolishes delay-period stability without degrading sensory perception \cite{devlin2018bert}.
\end{abstract}

\section{Introduction}
Persistent neural activity during memory delay periods has long been considered the cellular substrate of cognitive buffer maintenance \cite{vaswani2017attention}. However, recent theoretical models argue that dynamic attractor landscapes, rather than static firing rates, govern information retention in noisy cortical networks \cite{shimizu2003}.

\section{Results}
Electrophysiological recordings during a delayed non-match to sample task revealed robust attractor dynamics in dorsolateral prefrontal cortex \cite{nature2022}. Population vector trajectories formed stable low-dimensional manifolds during the retention interval \cite{radford2021learning}.

\bibliographystyle{nature}
\bibliography{references}
\end{document}
